% Chapter 5

\chapter{Conclusion} % Write in your own chapter title
\label{Chapter7}
\lhead{Chapter 7. \emph{Conclusion}} 
\section{Summary of Findings}
This project proposes a complete system that integrates the existing state-of-the-art temporal activity recognition model into a desktop application that leverages the existing camera infrastructure found in clinical settings to recognize patient activity and automate the logging and processing of these data. This ,in turn, assists  nurses in their tasks consequently improving patient outcomes. A prototype of the proposed system is also developed, which achieves 63.38\%  in proposal generation using the THUMOS'14 dataset. It forecasts  patient activities with an MAE of 33.1\% when using a general model and an MAE of 96.7\% when using a personalized model trained on the Toyota Smart Home dataset. The system also provides push notifications to nurses in the case of emergency alerts and provides doctors with realtime patient analytics along with Gemini's Generative AI capabilities to enhance data analysis.
\section{Benefits of the Project}
\subsection{Stakeholders}
\subsubsection{Elderly Patients}
The project aimed to improve  health by providing regular monitoring. The system detects and records unusual behaviors such as falls, sleeping patterns, eating habits, and vital signs, to ensure that any abnormalities are quickly addressed. As a result, elderly patients experience better treatment and  health outcomes.
\subsubsection{Nurses}
By automating the Nursing Care Process (NCP) , the project reduces nurses' workload. Nurses can plan their work, set priorities, and obtain real-time statistics on patient health by utilizing computer vision technologies and digitized nursing systems. Consequently, nurses are better able to care for their patients and experience less burnout.
\subsubsection{Doctors}
The project provides doctors with an extensive dashboard that provides a quick picture of a patient's activities, abnormalities, and vital signs. Doctors can make informed decisions and  offer more effective treatment plans when they have access to precise and current data.
\subsubsection{Caretakers}
The mobile application allows caregivers to access information about their patients, strengthening the relationship between the caretaker and  medical staff. This helps caretakers  build confidence in the quality of the care being delivered.
\subsection{Outputs Expected from the Project}
The outputs expected untill  the completion of the project are listed below:
\begin{enumerate}
    \item Provide healthcare providers with a comprehensive dashboard that displays real-time analytics of patients' health, including vital signs and abnormalities using infographics. Doctors will make informed decisions and treatments to, improve patient outcomes.
    \item The project seeks to enhance the standard of care and assistance provided to elderly individuals. The project will help to ensure timely treatment, reduce risks such as falls, and enhance the overall health of patients by keeping track of their movements and vital signs.
    \item By automating tasks and improving the nursing care workflow, this project will reduce the burden on nurses. Nurses will be able to prioritize their work more efficiently and provide patients with better care by providing real-time statistics and warnings.
    \item Caretakers  have access to the project's mobile application to receive updates and information on their patients. This will strengthen the relationship between the caretaker and medical-staff by establishing confidence and trust in the care provided.
    \item The project's application in healthcare institutions such as hospitals and nursing homes, will increase efficiency and reduce costs. Healthcare organizations will improve their reputations and attract more consumers by reducing nurse burnout and increasing patient care quality.
    \item This project will lay the foundation for future healthcare research and development. It explores how the latest technologies, such as computer vision and AI, can be used in nursing care operations. This new information will be useful to researchers and developers as they work to advance healthcare technology.
\end{enumerate}
\section{Recommendations for Future Research}
This project was a successful implementation of the fact that we could create a software application by leveraging  recent research  in the fields of Computer Vision and Generative AI. In this project, we were able to combine state of the art research to detect activities of daily living and generative AI, to create a major impact in the health care sector. A basic prototype has been built as a bachelor's degree project, but there are many modules that can be added in the future in an incremental phase. Some of the future work and research that can be conducted are as follows:

\begin{enumerate}
    \item To date, we have focused on recognizing activities of daily living in elderly patients. Undoubtedly, this will assist in the improvement of the lifestyles of elderly patients, but a more impactful idea would be to detect abnormal activities and respond to them effectively. For example, falling is a very common activity in elderly patients and a major cause of severe injuries such as fractures and joint dislocations. Therefore, adding a module to detect abnormal activities and alerting  care takers would surely be of more use.
    \item The main role of  care takers is to comfort  elderly patients by overseeing their daily living activities and providing them with assistance and medical care according to their needs. Under these circumstances, it is crucial to understand the behavior of  patients. Looking for patterns in them and how they evolve with time can help diagnose problems at an early stage and treat them in a timely manner. This process of behavior analysis can be automated by leveraging artificial intelligence to ensure that abnormalities in behaviors can be detected effectively.
    \item Integration of generative AI for Documentation and Decision support
    \item Integrating personalized models can revolutionized in the health care field. Each patient visiting a hospital has his own needs and way of living. Working with generic models can be a way to enter the market of automating the health care sector, but in the longer run, moving to personalized models will be the game changer.
    \item Sensors can be integrated with  applications to automate  nursing activities. For example, one of the most time consuming activities of a nurse is to manually record the vital signs of a patient. We can leverage state of the art wrist watches and bands that have the capability to measure  blood pressure, heart rate, and oxygen levels. Integrating these sensors will empower  nurses to have more time to help patients emotionally.
    \item Face recognition and object tracking can be integrated into an application to deal with situations in which multiple people are in a frame. This will help the model  exclude other persons in the frame and focus on detecting the activities of the patient.
    \item Currently, a single camera is used to detect patient activities of the patient. In a future work, we will integrate multiple cameras in a room or environment that can work in synchronization to monitor the patients in their complete surrounding without bounding them at a certain location in the room.
    
    

\end{enumerate}
% Write in your own chapter title to set the page header

